\section{Related Works}~\label{sec:related-works}

Research on 3D descriptors for place recognition and SLAM has evolved from classical handcrafted methods to learning-based and multi-modal fusion approaches. Early techniques relied on manually designed geometric features extracted from point clouds, either locally or globally, providing efficient but limited representations. More recently, deep learning methods operating solely on point clouds have demonstrated superior discriminative power by automatically learning invariant and task-optimized features. However, since LiDAR geometry alone may be insufficient in perceptually ambiguous or long-term changing environments, there is growing interest in learning descriptors from fused sensory information, particularly LiDAR–camera combinations, to exploit complementary geometric and visual cues. Motivated by this, this section reviews first the analyzes of classical geometric descriptors, then learning-based methods using only point clouds, and finally approaches based on LiDAR–camera fusion, highlighting the importance of multi-modal learned representations for robust long-term localization.

\subsection{LiDAR-Based Handcrafted Local and Global Descriptors}~\label{subsec:lidar-only-local-global-descriptors}

A LiDAR-based local descriptor is a compact representation extracted from the geometric neighborhood of a 3D point in a LiDAR point cloud. It encodes local spatial structure, shape distribution, and possibly reflectance information to facilitate correspondence estimation, feature matching, or higher-level scene understanding tasks.

Table~\ref{tab:lidar-only-local-descriptors} summarizes key handcrafted local and global descriptors for LiDAR point clouds, highlighting their approaches, invariance properties, strengths, and limitations. Handcrafted methods such as Spin Images \cite{Johnson1999} and LiDAR intensity-based descriptors \cite{HungarConstanze2019} rely on geometric and intensity gradients to capture local structure, offering robustness to illumination changes but often at high computational cost. Learning-based approaches like USIP \cite{Li2019} and PPF-AE \cite{deng2018} leverage unsupervised learning to identify stable keypoints and encode geometric relationships, providing improved repeatability and robustness but with increased complexity. Overall, while handcrafted descriptors laid the foundation for local feature extraction in 3D point clouds, learning-based methods have shown superior performance by automatically optimizing representations for specific tasks, albeit with trade-offs in computational efficiency.

While LiDAR-based global descriptor is a holistic representation computed from a full 3D LiDAR scan, encoding global geometric layout and structural patterns of the environment into a compact embedding suitable for place recognition and loop closure detection.

These descriptors range from geometric and projection-based methods to segmentation-driven approaches, each with unique advantages and challenges. For instance, semantic-enriched descriptors \cite{Fan2021, Liao2022} offer improved discriminability but depend heavily on accurate segmentation, while geometric descriptors \cite{Cao2018, Zhang2022} are effective in structured environments but may struggle in unstructured scenes. Projection-based methods \cite{Wang2020, Wang2023} provide compact representations but can lose fine 3D details. Overall, while these handcrafted global descriptors have been foundational in LiDAR-based place recognition and localization, they often face limitations in dynamic or changing environments, motivating the shift towards learning-based and multi-modal approaches for enhanced robustness.

%LiDAR-only semantic segmentation methods aim to assign a semantic label to each 3D point in a point cloud using exclusively geometric information captured by the sensor. These approaches rely on the spatial structure, depth measurements, an local geometric properties of the environment without incorporating visual texture cues from cameras. Broadly, LiDAR-based segmentation techniques can be categorized into classical geometric descriptors based on points cloud, voxel, image methods.  

%\subsubsection{Point cloud-based methods} \label{subsubsec:point-cloud-method}

%Early methods such as \cite{qi2017pointnetdeeplearningpoint} and \cite{qi2017pointnet++} employ Multi-Layer Perceptrons (MLPs) combined with KNN-based hierarchical grouping to learn local geometric structures from raw point sets. These approaches rely on hierarchical aggregation mechanisms to progressively capture neighborhood information. Similarly, \cite{qian2022pointnextrevisitingpointnetimproved} extends MLP-based hierarchical architectures by incorporating efficient MLP stacking and explicit KNN neighbor search to improve scalability and training stability.
%More advanced point-based convolutional models are introduced in \cite{thomas2024kpconvxmodernizingkernelpoint}, where kernel point convolution combined with attention mechanisms enhances local feature modeling. In \cite{hu2020randlanetefficientsemanticsegmentation}, attention is further integrated into point-based MLP networks with detection-aware optimization, incorporating random sampling and local feature aggregation strategies evaluated on SemanticKITTI. Additionally, \cite{yan20222dpass2dpriorsassisted} proposes a 3D point network enhanced through knowledge distillation from 2D priors, enabling semantic transfer from image-based representations into 3D space.
%Attention mechanisms and global context modeling become central in later works. The self-attention framework in \cite{lai2023sphericaltransformerlidarbased3d} adopts a transformer-style architecture operating on spherical-coordinate point representations, using radial window self-attention to capture global structural dependencies. In contrast, \cite{sanchez20253dlabelpropgeometricdrivendomaingeneralization} introduces a geometric label propagation strategy based on point alignment and ICP propagation, focusing on geometric optimization rather than end-to-end supervised learning.
%Transformer-based architectures dominate the most recent developments. Works \cite{zhao2021pointtransformer, NEURIPS2022_d78ece66, wu2024pointtransformerv3simpler} propose self-attention-driven models that operate directly on raw point sets. Specifically, [126] incorporates positional encoding into self-attention modules, [127] introduces efficient neighborhood attention with residual context blocks for end-to-end optimization, and [128] proposes a multi-scale global–local transformer fusion framework with computational optimizations to improve robustness in large-scale outdoor scenes.

%\subsubsection{Voxel-based methods} \label{subsubsec:voxel-based}

%\cite{Maturana2015} proposes a model architecture to address this problem by integrating a volumetric occupancy grid representation with a supervised 3D Convolutional Neural Network (3D CNN). Their approach guarantees invariance to rotation and resolution by employing fewer CNN parameters compared to other existing models, while still achieving reliable performance in object detection tasks.
%MinkLoc3D \cite{komorowski2020minkloc3dpointcloudbased} introduces a framework based on sparse 3D convolutions combined with a Feature Pyramid Network (FPN) backbone and GeM pooling to generate global descriptors. Similarity is computed using L2 distance between embeddings. MinkLoc3D-SI \cite{Zywanowski_2022} extends this approach by incorporating spherical coordinates and intensity information into the sparse convolution pipeline. It employs a triplet loss with explicit L2 distance between descriptors.
%EgoNN \cite{komorowski2021egonnegocentricneuralnetwork} leverages an FPN backbone enhanced with Efficient Channel Attention (ECA) modules to improve feature representation. It adopts global descriptor retrieval for place recognition, emphasizing improved contextual feature aggregation. Similarly, NDT-Transformer \cite{ZhouZhicheng2021} combines voxel-level statistical representations based on the Normal Distribution Transform (NDT) with transformer encoders to produce global descriptors,
%TransLoc3D \cite{xu2022transloc3dpointcloud} integrates NDT representations with an Adaptive Receptive Field Module (ARFM) and transformer architecture, followed by NetVLAD aggregation for global descriptor extraction. It uses triplet margin loss with L2 distance between descriptors, demonstrating robust performance on Oxford RobotCar. LCDNet \cite{cattaneo2022lcdnetdeeploopclosure} adopts a PV-RCNN backbone for feature extraction and employs NetVLAD for global descriptor aggregation, using L2 or cosine similarity metrics. It is evaluated on KITTI, KITTI-360, and Freiburg datasets.
%SVT-Net \cite{fan2021svtnetsuperlightweightsparse} introduces a Sparse Voxel Transformer (SVT) that captures both local and long-range structural dependencies before global descriptor retrieval, highlighting the importance of modeling extended spatial context. Finally, OverlapNetVLAD \cite{Fu2023ACP} employs a BEV-based feature extractor (BEVNet) combined with NetVLAD to produce descriptors typically compared via L2 distance, showing effectiveness on KITTI (2023).

%\subsubsection{Image-based methods} \label{subsubsec:image-based-methods}

%KPConv \cite{sanchez20253dlabelpropgeometricdrivendomaingeneralization} directly processes raw point clouds using Kernel Point Convolution, emphasizing geometric domain adaptation to improve generalization on SemanticKITTI. BEV-based convolutional fusion \cite{qiu2024pcbevefficientpolarcartesianbev} projects LiDAR data into polar and Cartesian Bird’s-Eye View representations and employs dual-branch CNNs with joint spatial–semantic optimization, evaluated on nuScenes.
%Hybrid approaches such as PointConv + SparseConv3D \cite{park2022pcscnetfast3dsemantic} operate on voxelized 3D grids, combining sparse 3D convolutions with point-based refinement in an end-to-end supervised framework. Harmonic Dense Convolution \cite{razani2021litehdseglidarsemanticsegmentation} also relies on voxel representations but introduces harmonic dense filters to improve computational efficiency. 
%Projection-based methods include 2D CNNs over spherical projections \cite{cortinhal2024salsanextfastuncertaintyawaresemantic}, which incorporate residual context and Bayesian uncertainty modeling, and 2D CNNs over polar range images \cite{nowruzi2021polarnetaccelerateddeepopen}, which apply standard supervised training on structured range-view representations.

\subsection{Learning-Based LiDAR Local and Global Descriptors}~\label{subsec:learning-based-local-global-descriptors}

Learning-based local descriptors are feature representations automatically learned from data using neural networks, encoding the geometric and/or semantic structure of local neighborhoods within a 2D or 3D space. This strategy is divide into 5 stages, as depicted in Table~\ref{tab:learning-based-lidar}: 1) point-based methods, which directly operate on raw point clouds; 2) voxel-based methods, which discretize the point cloud into a voxel grid and apply 3D convolutions; 3) segment-based methods, which extract local segments from the point cloud and learn descriptors based on their geometric and semantic properties; 4) projection-based methods, which project the 3D point cloud onto 2D planes; and 5) graph-based methods, which represent the point cloud as a graph structure and learn descriptors based on node and edge features.

Descriptor evolution across the references reflects increasing abstraction and contextual modeling. Early point-based descriptors such as PointNetVLAD~\cite{uy2018pointnetvladdeeppointcloud} and its extensions~\cite{liu2019lpdnet3dpointcloud, du2020dh3ddeephierarchical3d, HuiLe2021}, rely on NetVLAD aggregation for global feature pooling. 

\renewcommand{\arraystretch}{1.2}
\begin{table*}[!t]
\caption{\scriptsize\MakeUppercase{Learning-Based LiDAR Local and Global Descriptors}}
\label{tab:learning-based-lidar}
\centering
\begin{tabular}{|c|c|c|c|c|}
\hline
\rowcolor[HTML]{EFEFEF} 
\textbf{Ref.} & \textbf{Class} & \textbf{Descriptor} & \textbf{Feature} & \textbf{Similarity} \\ \hline
\cite{uy2018pointnetvladdeeppointcloud} &  & PointNet & NetVLAD & L2 / Cosine \\ \cline{1-1} \cline{3-5} 
\cite{liu2019lpdnet3dpointcloud} &  & GNN, DG-CNN & NetVLAD & Cosine \\ \cline{1-1} \cline{3-5} 
\cite{du2020dh3ddeephierarchical3d} &  & PointNet + FlexConv & NetVLAD, SE block & Cosine \\ \cline{1-1} \cline{3-5} 
\cite{HuiLe2021} &  & \begin{tabular}[c]{@{}c@{}}Pyramid Point \\ Transformer\end{tabular} & Pyramid VLAD & L2 / Cosine \\ \cline{1-1} \cline{3-5} 
\cite{Vidanapathirana_2022} & \multirow{-5}{*}{Point-based} & U-Net & ePN, 2nd-order pooling & L2 \\ \hline
\cite{komorowski2020minkloc3dpointcloudbased} &  & MinkLoc3D & Sparse 3D convolutions, FPN, GeM & L2 \\ \cline{1-1} \cline{3-5} 
\cite{Zywanowski_2022} &  & MinkLoc3D-SI & \begin{tabular}[c]{@{}c@{}}Sparse convolutions - spherical \\ coordinates - intensity; FPN; GeM\end{tabular} & L2 \\ \cline{1-1} \cline{3-5} 
\cite{komorowski2021egonnegocentricneuralnetwork} &  & EgoNN & FPN backbone; ECA attention & \begin{tabular}[c]{@{}c@{}}Uses global \\ descriptor retrieval\end{tabular} \\ \cline{1-1} \cline{3-5} 
\cite{ZhouZhicheng2021} &  & NDT-Transformer & \begin{tabular}[c]{@{}c@{}}Voxel-level statistical representation \\ (Normal Distribution Transform); \\ Transformer encoder;\end{tabular} & \begin{tabular}[c]{@{}c@{}}Uses descriptor \\ retrieval\end{tabular} \\ \cline{1-1} \cline{3-5} 
\cite{xu2022transloc3dpointcloud} & \multirow{-5}{*}{Voxel-based} & TransLoc3D & \begin{tabular}[c]{@{}c@{}}NDT representation; Adaptive \\ Receptive Field Module (ARFM); \\ Transformer; NetVLAD\end{tabular} & L2 \\ \hline
\cite{LiLin2021} &  & SA-LOAM & \begin{tabular}[c]{@{}c@{}}Edge and planar features (LOAM style); \\ spatial-appearance cues\end{tabular} & ICP-based \\ \cline{1-1} \cline{3-5} 
\cite{li2021sscsemanticscancontext} &  & SSC & \begin{tabular}[c]{@{}c@{}}PCA-based shape descriptors \\ from segmented point clouds\end{tabular} & \begin{tabular}[c]{@{}c@{}}L2 between \\ PCA descriptors\end{tabular} \\ \cline{1-1} \cline{3-5} 
\cite{fan2021semantic} &  & PCA-SSC & \begin{tabular}[c]{@{}c@{}}Improved PCA descriptors with \\ stronger geometric invariance\end{tabular} & L2 distance \\ \cline{1-1} \cline{3-5} 
\cite{yin2021psematchviewpointfreeplacerecognition} &  & PSE-Match & \begin{tabular}[c]{@{}c@{}}Probabilistic Shape Encoding \\ (mixture-model representation)\end{tabular} & \begin{tabular}[c]{@{}c@{}}Probability-based \\ similarity measure\end{tabular} \\ \cline{1-1} \cline{3-5} 
\cite{TaoSong2022} & \multirow{-5}{*}{Segment-based} & Song et al. & \begin{tabular}[c]{@{}c@{}}Hybrid descriptor combining geometry, \\ intensity, and local topology\end{tabular} & L2 or cosine \\ \hline
\cite{Barros_2022} &  & AttDLNet & \begin{tabular}[c]{@{}c@{}}Attention-based CNN encoder on projection \\ images; multi-scale feature extraction\end{tabular} & L2 \\ \cline{1-1} \cline{3-5} 
\cite{YinFuying2021} &  & SphereVLAD & \begin{tabular}[c]{@{}c@{}}Spherical projection + deep CNN features\\  + VLAD aggregation\end{tabular} & L2 / cosine distance \\ \cline{1-1} \cline{3-5} 
\cite{zhao2022spherevladattentionbasedsignalenhancedviewpoint} &  & SphereVLAD++ & \begin{tabular}[c]{@{}c@{}}Improved SphereVLAD with multi-scale \\ projections and enhanced CNN backbone\end{tabular} & L2 distance \\ \cline{1-1} \cline{3-5} 
\cite{lu2022deepringlearningrototranslationinvariant} &  & DeepRING & \begin{tabular}[c]{@{}c@{}}Circular (ring-based) LiDAR projection; \\ CNN encoder for rotation-equivariant features\end{tabular} & \begin{tabular}[c]{@{}c@{}}Correlation-based \\ similarity\end{tabular} \\ \cline{1-1} \cline{3-5} 
\cite{luo2023bevplacelearninglidarbasedplace} & \multirow{-5}{*}{Projection-based} & BEVPlace & \begin{tabular}[c]{@{}c@{}}Bird’s-eye-view (BEV) projection; CNN \\ encoder with BEV-specific geometric priors\end{tabular} & L2 distance \\ \hline
\cite{kong2020semanticgraphbasedplace} &  & SGPR & \begin{tabular}[c]{@{}c@{}}Builds semantic graphs from \\ segmented point clouds\end{tabular} & \begin{tabular}[c]{@{}c@{}}Graph matching using \\ structural similarity metrics\end{tabular} \\ \cline{1-1} \cline{3-5} 
\cite{Zhu2020GOSMatchGM} &  & GOS Match & \begin{tabular}[c]{@{}c@{}}Graph-of-semantics representation where \\ each node = semantic instance \\ and edges = geometric/topological relations\end{tabular} & \begin{tabular}[c]{@{}c@{}}Graph similarity scoring \\ with node - edge\end{tabular} \\ \cline{1-1} \cline{3-5} 
\cite{GongYansong2021} &  & Gong et al. & \begin{tabular}[c]{@{}c@{}}Two-tier representation: local spatial relation \\ graph + global spatial topology; \\ semantic and geometric cues\end{tabular} & \begin{tabular}[c]{@{}c@{}}Coarse-to-fine graph \\ matching;\end{tabular} \\ \cline{1-1} \cline{3-5} 
\cite{pramatarov2021a} &  & BoxGraph & \begin{tabular}[c]{@{}c@{}}Semantic bounding box graph generated from \\ object detectors; boxes as nodes, \\ box relations as edges\end{tabular} & \begin{tabular}[c]{@{}c@{}}Graph matching with \\ similarity metrics\end{tabular} \\ \cline{1-1} \cline{3-5} 
\cite{cui2021dscdeepscancontext} & \multirow{-5}{*}{Graph-based} & \begin{tabular}[c]{@{}c@{}}Deep Scan\\ Context\end{tabular} & \begin{tabular}[c]{@{}c@{}}Learned high-dimensional scan context descriptor; \\ pseudo-cylindrical projection + neural embedding\end{tabular} & L2 \\ \hline
\end{tabular}
\end{table*}
\renewcommand{\arraystretch}{1.2}

Sparse convolution-based descriptors like MinkLoc3D~\cite{komorowski2020minkloc3dpointcloudbased} and MinkLoc3D-SI \cite{Zywanowski_2022} introduce hierarchical 3D feature encoding. EgoNN \cite{komorowski2021egonnegocentricneuralnetwork} integrates attention mechanisms for enhanced global retrieval.
Voxel-based descriptors such as NDT-Transformer~\cite{ZhouZhicheng2021} and TransLoc3D~\cite{xu2022transloc3dpointcloud} employ statistical encodings and transformer architectures to capture long-range spatial dependencies.
Segment-based descriptors like SSC~\cite{li2021sscsemanticscancontext}, PCA-SSC~\cite{fan2021semantic}, and PSE-Match~\cite{yin2021psematchviewpointfreeplacerecognition} focus on geometric invariance and probabilistic similarity modeling. Song et al.~\cite{TaoSong2022} propose hybrid descriptors combining geometry, intensity, and topology.
Projection-based descriptors including ADLNet~\cite{Barros_2022}, SphereVLAD~\cite{YinFuying2021}, SphereVLAD++~\cite{zhao2022spherevladattentionbasedsignalenhancedviewpoint}, DeepRING~\cite{lu2022deepringlearningrototranslationinvariant}, and BEVPlace~\cite{luo2023bevplacelearninglidarbasedplace} utilize CNN encoders on structured projections, incorporating multi-scale aggregation and rotation-equivariant mechanisms.
Graph-based descriptors introduce a structural paradigm shift. SGPR~\cite{kong2020semanticgraphbasedplace} and GNN-based matching methods~\cite{Zhu2020GOSMatchGM} encode semantic graphs and perform graph similarity learning. Cong et al.~\cite{GongYansong2021} integrate two-branch networks combining local spatial relations and graph-level topology. Box-Graph~\cite{pramatarov2021a} and DGC Scan Context~\cite{cui2021dscdeepscancontext} further exploit structural descriptors and graph matching metrics for robust place recognition.
Overall, descriptor development evolves from global feature aggregation toward relational graph modeling and topology-aware embeddings. \\

Feature representations vary significantly across paradigms. Point-based works such as U-Net~\cite{Vidanapathirana_2022} and MinkLoc3D~\cite{komorowski2020minkloc3dpointcloudbased} leverage sparse 3D convolutions, second-order pooling, and geometric descriptors. These methods focus on hierarchical geometric encoding.
Voxel-based features in NDT-Transformer~\cite{ZhouZhicheng2021} rely on statistical distribution modeling (Normal Distribution Transform), while TransLoc3D~\cite{xu2022transloc3dpointcloud} introduces adaptive receptive field modules and NetVLAD aggregation for enhanced contextual capture.
Segment-based approaches such as SSC~\cite{li2021sscsemanticscancontext} and PCA-SSC~\cite{fan2021semantic} employ PCA-derived geometric features, whereas PSE-Match~\cite{yin2021psematchviewpointfreeplacerecognition} models shapes probabilistically. Song et al.~\cite{TaoSong2022} integrate hybrid feature combinations including geometry and intensity.
Projection-based methods like SphereVLAD~\cite{YinFuying2021}, DeepRING~\cite{lu2022deepringlearningrototranslationinvariant}, and BEVPlace~\cite{luo2023bevplacelearninglidarbasedplace} extract features via CNN backbones applied to spherical or BEV projections, often incorporating rotation invariance and multi-scale representations.
Graph-based features explicitly encode nodes (e.g., semantic objects, segments, scan contexts) and edges (spatial relationships, topological adjacency). SGPR~\cite{kong2020semanticgraphbasedplace} builds semantic graphs, while [59] and [60] incorporate attention mechanisms and relational reasoning. Box-Graph~\cite{pramatarov2021a} and DGC Scan Context~\cite{cui2021dscdeepscancontext} emphasize structural graph descriptors and pose-invariant relational encoding.
Thus, feature design increasingly transitions from local geometric encoding toward structured relational representations. \\

Similarity computation strategies reflect methodological differences across paradigms. Most deep descriptor-based approaches, including~\cite{uy2018pointnetvladdeeppointcloud, HuiLe2021, Vidanapathirana_2022, komorowski2020minkloc3dpointcloudbased, Zywanowski_2022, xu2022transloc3dpointcloud, li2021sscsemanticscancontext, fan2021semantic, TaoSong2022, Barros_2022,YinFuying2021, zhao2022spherevladattentionbasedsignalenhancedviewpoint, luo2023bevplacelearninglidarbasedplace, cui2021dscdeepscancontext}, rely on L2 or cosine similarity in embedding space for efficient large-scale retrieval.
NetVLAD-based methods~\cite{uy2018pointnetvladdeeppointcloud,du2020dh3ddeephierarchical3d}, primarily use cosine similarity. SSC-based methods~\cite{li2021sscsemanticscancontext, fan2021semantic}, use L2 distance on PCA-derived descriptors. PSE-Match~\cite{yin2021psematchviewpointfreeplacerecognition} employs a probabilistic similarity measure. SA-LOAM~\cite{TaoSong2022} is distinct in using ICP-based alignment.
DeepRING~\cite{lu2022deepringlearningrototranslationinvariant} introduces correlation-based similarity for rotation robustness. Voxel-based NDT-Transformer~\cite{ZhouZhicheng2021} emphasizes descriptor retrieval frameworks.
Graph-based approaches adopt more complex structural matching strategies. SGPR \cite{kong2020semanticgraphbasedplace} uses graph matching and structural similarity scoring. GNN-based matching~\cite{Zhu2020GOSMatchGM} integrates graph similarity learning with node and edge-level reasoning. Cong et al.~\cite{GongYansong2021} apply contrastive graph matching. Box-Graph~\cite{pramatarov2021a} and DGC Scan Context~\cite{cui2021dscdeepscancontext} employ graph similarity metrics and L2 distance over relational descriptors.
Overall, similarity computation evolves from simple metric comparisons toward structured graph matching and relational reasoning, trading computational simplicity for structural robustness.

\subsection{Camera-LiDAR fusion methods} \label{subsec:camera-lidar-methods}

\begin{table*}[!t]
\caption{\scriptsize\MakeUppercase{Camera-LiDAR-fusion descriptors}}
\label{tab:camera-lidar-fusion-descriptors}
\centering
\begin{tabular}{|c|c|c|c|c|c|}
\hline
\rowcolor[HTML]{EFEFEF} 
\textbf{Ref.} & \textbf{Descriptor Type} & \textbf{Primary task} & \textbf{Multi modal Strategy} & \textbf{Strengths} & \textbf{Limitation} \\ \hline
\cite{komorowski2021minkloclidarmonocularimage} &  &  & \begin{tabular}[c]{@{}c@{}}Late fusion of LiDAR + \\ image embeddings\end{tabular} & \begin{tabular}[c]{@{}c@{}}Combines geometric and visual cues; \\ strong generalization\end{tabular} & \begin{tabular}[c]{@{}c@{}}Requires synchronized L-C; \\ moderate computational cost\end{tabular} \\ \cline{1-1} \cline{4-6} 
\cite{pan2021coralcoloredstructuralrepresentation} &  &  & \begin{tabular}[c]{@{}c@{}}Colored structural + \\ visual fusion\end{tabular} & \begin{tabular}[c]{@{}c@{}}Integrates appearance and structure; \\ robust place recognition\end{tabular} & \begin{tabular}[c]{@{}c@{}}Slightly lower recall than top-\\ performing multimodal methods\end{tabular} \\ \cline{1-1} \cline{4-6} 
\cite{Xu2025} &  &  & \begin{tabular}[c]{@{}c@{}}Cross-modal alignment - \\ orientation voting\end{tabular} & \begin{tabular}[c]{@{}c@{}}Orientation robustness via voting; \\ cross-modal alignment\end{tabular} & \begin{tabular}[c]{@{}c@{}}Sensor calibration complexity; \\ computational overhead\end{tabular} \\ \cline{1-1} \cline{4-6} 
\cite{Zhou2023} &  &  & \begin{tabular}[c]{@{}c@{}}Multi-scale LiDAR–camera \\ attention fusion\end{tabular} & \begin{tabular}[c]{@{}c@{}}Captures geometric and appearance \\ cues; robust under viewpoint change\end{tabular} & \begin{tabular}[c]{@{}c@{}}Accurate cross-modal calibration; \\ attention increases computational cost\end{tabular} \\ \cline{1-1} \cline{4-6} 
\cite{LIANG2025408} & \multirow{-5}{*}{\begin{tabular}[c]{@{}c@{}}Global \\ multimodal \\ descriptor\end{tabular}} & \multirow{-5}{*}{\begin{tabular}[c]{@{}c@{}}Place \\ Recognition\end{tabular}} & \begin{tabular}[c]{@{}c@{}}Bi-directional attentional \\ fusion\end{tabular} & \begin{tabular}[c]{@{}c@{}}Attentional refinement; resilient \\ to ambiguous scenes\end{tabular} & \begin{tabular}[c]{@{}c@{}}Higher computational cost due \\ to attention mechanism\end{tabular} \\ \hline
\cite{Li2023} &  &  & \begin{tabular}[c]{@{}c@{}}Multi-modal driving \\ fusion (RGB + LiDAR)\end{tabular} & \begin{tabular}[c]{@{}c@{}}Accurate multimodal segmentation; \\ integrates camera + LiDAR\end{tabular} & \begin{tabular}[c]{@{}c@{}}Computationally heavy; dataset-\\ specific performance\end{tabular} \\ \cline{1-1} \cline{4-6} 
\cite{Zhang2023} &  &  & \begin{tabular}[c]{@{}c@{}}Cross-modal transformer \\ fusion (RGB-X)\end{tabular} & \begin{tabular}[c]{@{}c@{}}Strong cross-modal feature learning; \\ improved semantic accuracy\end{tabular} & \begin{tabular}[c]{@{}c@{}}Transformer complexity; \\ high GPU memory usage\end{tabular} \\ \cline{1-1} \cline{4-6} 
\cite{hou2022pointtovoxelknowledgedistillationlidar} &  & \multirow{-3}{*}{\begin{tabular}[c]{@{}c@{}}Semantic \\ Segmentation\end{tabular}} & \begin{tabular}[c]{@{}c@{}}LiDAR + teacher features \\ (cross-modal distillation)\end{tabular} & \begin{tabular}[c]{@{}c@{}}Unsupervised knowledge distillation; \\ improves LiDAR descriptors\end{tabular} & \begin{tabular}[c]{@{}c@{}}Dependent on teacher model; \\ added training complexity\end{tabular} \\ \cline{1-1} \cline{3-6} 
\cite{2023PatRe.13709299H} &  &  & \begin{tabular}[c]{@{}c@{}}Multi-modal semantic \\ adaptation\end{tabular} & \begin{tabular}[c]{@{}c@{}}Handles domain shift; adapts \\ features without supervision\end{tabular} & \begin{tabular}[c]{@{}c@{}}Limited semantic precision compared \\ to fully supervised methods\end{tabular} \\ \cline{1-1} \cline{4-6} 
\cite{Cao2024} &  & \multirow{-2}{*}{\begin{tabular}[c]{@{}c@{}}Unsupervised \\ Domain \\ Adaptation\end{tabular}} & \begin{tabular}[c]{@{}c@{}}Multi-modal prior-guided \\ cross-domain adaptation\end{tabular} & \begin{tabular}[c]{@{}c@{}}Leverages image priors to reduce \\ domain gap; improves target \\ generalization\end{tabular} & \begin{tabular}[c]{@{}c@{}}Relies on quality of pseudo-labels; \\ sensitive to domain discrepancy\end{tabular} \\ \cline{1-1} \cline{3-6} 
\cite{LuoYang2025} &  &  & \begin{tabular}[c]{@{}c@{}}Cross-modal semantic \\ attention\end{tabular} & \begin{tabular}[c]{@{}c@{}}Focused semantic fusion; \\ enhances point-level semantics\end{tabular} & \begin{tabular}[c]{@{}c@{}}Requires accurate alignment \\ between sensors\end{tabular} \\ \cline{1-1} \cline{4-6} 
\cite{Lin2025} &  &  & Vision assistance for LiDAR & \begin{tabular}[c]{@{}c@{}}Leverages visual context to \\ improve LiDAR features\end{tabular} & \begin{tabular}[c]{@{}c@{}}Needs dense camera coverage; \\ limited generalization if camera fails\end{tabular} \\ \cline{1-1} \cline{4-6} 
\cite{Snchez-Garca2025} &  &  & Range + image fusion & \begin{tabular}[c]{@{}c@{}}Real-time performance; \\ robust multimodal segmentation\end{tabular} & \begin{tabular}[c]{@{}c@{}}Requires synchronized RGB; \\ may struggle with sparse LiDAR\end{tabular} \\ \cline{1-1} \cline{4-6} 
\cite{Lamichhane2025} &  &  & \begin{tabular}[c]{@{}c@{}}Depth-aware transformer \\ fusion\end{tabular} & \begin{tabular}[c]{@{}c@{}}Attention-based multimodal \\ context; strong semantic consistency\end{tabular} & \begin{tabular}[c]{@{}c@{}}Transformer complexity; \\ high memory usage\end{tabular} \\ \cline{1-1} \cline{4-6} 
\cite{Zhang2025} &  &  & LiDAR + fisheye fusion & \begin{tabular}[c]{@{}c@{}}Wide field-of-view coverage; \\ captures rich context\end{tabular} & \begin{tabular}[c]{@{}c@{}}Needs fisheye camera; \\ calibration-sensitive\end{tabular} \\ \cline{1-1} \cline{4-6} 
\cite{Tan2024} &  &  & Efficient perception fusion & \begin{tabular}[c]{@{}c@{}}Efficient fusion; perception-\\ aware features\end{tabular} & \begin{tabular}[c]{@{}c@{}}Focused on efficiency, may trade \\ off accuracy in complex scenarios\end{tabular} \\ \cline{1-1} \cline{4-6} 
\cite{Ni2024Robust3S} &  &  & \begin{tabular}[c]{@{}c@{}}Collaborative multimodal \\ learning\end{tabular} & \begin{tabular}[c]{@{}c@{}}Robust semantic feature \\ learning across modalities\end{tabular} & \begin{tabular}[c]{@{}c@{}}Higher training complexity; \\ needs aligned multimodal data\end{tabular} \\ \cline{1-1} \cline{4-6} 
\cite{Duan2024} &  &  & \begin{tabular}[c]{@{}c@{}}Camera–LiDAR cross-attention \\ with fast neighbor aggregation\end{tabular} & \begin{tabular}[c]{@{}c@{}}Efficient local-global feature \\ interaction; enhanced boundary \\ precision\end{tabular} & \begin{tabular}[c]{@{}c@{}}Attention and neighbor aggregation \\ increase memory usage\end{tabular} \\ \cline{4-6} 
\cite{Zhao2024} &  &  & \begin{tabular}[c]{@{}c@{}}Early–mid LiDAR–image \\ feature fusion\end{tabular} & \begin{tabular}[c]{@{}c@{}}Improves discriminative power of \\ LiDAR features using visual cues\end{tabular} & \begin{tabular}[c]{@{}c@{}}Dependent on synchronized RGB; \\ performance drops under poor lighting\end{tabular} \\ \cline{1-1} \cline{4-6} 
\cite{Du2023} &  &  & \begin{tabular}[c]{@{}c@{}}Feature-level RGB–LiDAR \\ fusion with hierarchical encoder\end{tabular} & \begin{tabular}[c]{@{}c@{}}Enhanced small-object segmentation; \\ better contextual modeling\end{tabular} & \begin{tabular}[c]{@{}c@{}}Dataset-dependent tuning; moderate \\ computational overhead\end{tabular} \\ \cline{1-1} \cline{4-6} 
\cite{Maiti_2023_CVPR} &  &  & \begin{tabular}[c]{@{}c@{}}Transformer-based multimodal \\ fusion\end{tabular} & \begin{tabular}[c]{@{}c@{}}Strong long-range cross-modal int.; \\ improved semantic consistency\end{tabular} & \begin{tabular}[c]{@{}c@{}}High GPU memory consumption; \\ complex training\end{tabular} \\ \cline{1-1} \cline{4-6} 
\cite{Wang2024} &  &  & \begin{tabular}[c]{@{}c@{}}Unified multimodal encoder \\ (optical + auxiliary modalities)\end{tabular} & \begin{tabular}[c]{@{}c@{}}Cost-effective design; scalable \\ multimodal integration\end{tabular} & \begin{tabular}[c]{@{}c@{}}Designed for remote sensing; limited \\ direct transfer to autonomous driving\end{tabular} \\ \cline{1-1} \cline{4-6} 
\cite{Ma2024} &  & \multirow{-13}{*}{\begin{tabular}[c]{@{}c@{}}Semantic \\ Segmentation\end{tabular}} & \begin{tabular}[c]{@{}c@{}}Hierarchical transformer-based \\ multimodal fusion\end{tabular} & \begin{tabular}[c]{@{}c@{}}Multi-scale contextual modeling; \\ strong global reasoning\end{tabular} & \begin{tabular}[c]{@{}c@{}}Computationally intensive; \\ large training data requirement\end{tabular} \\ \cline{1-1} \cline{3-6} 
\cite{Lu2024} &  & \begin{tabular}[c]{@{}c@{}}Semantic Scene \\ Completion\end{tabular} & \begin{tabular}[c]{@{}c@{}}Continuous fusion in \\ voxelized 3D grid\end{tabular} & \begin{tabular}[c]{@{}c@{}}Dense 3D context modeling; \\ consistent volumetric representation\end{tabular} & \begin{tabular}[c]{@{}c@{}}High memory usage due to voxel.;\\  computationally heavy\end{tabular} \\ \cline{1-1} \cline{3-6} 
\cite{Zhao2022} & \multirow{-20}{*}{\begin{tabular}[c]{@{}c@{}}Semantic \\ feature \\ descriptor\end{tabular}} & \begin{tabular}[c]{@{}c@{}}3D Vehicle \\ Detection\end{tabular} & \begin{tabular}[c]{@{}c@{}}Semantic-guided camera–LiDAR \\ feature augmentation\end{tabular} & \begin{tabular}[c]{@{}c@{}}Enhances object-aware fusion; \\ improves detection robustness\end{tabular} & \begin{tabular}[c]{@{}c@{}}Reliable semantic supervision; \\ added training complexity\end{tabular} \\ \hline
\end{tabular}
\end{table*}

The following analysis examines recent multimodal descriptor approaches by organizing them according to descriptor type, primary task, fusion strategy, strengths, and limitations. By structuring the comparison column-wise, this discussion highlights the evolution from global multimodal embeddings for place recognition toward semantically rich feature representations for dense scene understanding, while also identifying common design trends, performance advantages, and practical challenges associated with multimodal perception systems.
The Table~\ref{tab:camera-lidar-fusion-descriptors} distinguishes between Global multimodal descriptors (\cite{komorowski2021minkloclidarmonocularimage}–\cite{LIANG2025408}) and Semantic feature descriptors (\cite{Snchez-Garca2025,Lin2025},\cite{Li2023}-\cite{Zhao2022}), reflecting two fundamentally different design objectives. Global multimodal descriptors are primarily designed for place recognition, where LiDAR and visual data are fused to generate compact embeddings suitable for large-scale retrieval. These approaches emphasize discriminative global representation, cross-modal consistency, and robustness to viewpoint or illumination changes. In contrast, semantic feature descriptors target dense scene understanding tasks such as semantic segmentation, domain adaptation, scene completion, and object detection. Rather than producing a single global vector, these methods learn structured, context-rich feature maps that enable fine-grained prediction. This categorization highlights the methodological shift from retrieval-oriented embeddings to dense multimodal perception frameworks.

The section of primary tasks is boarded from global localization to comprehensive scene interpretation. Works~\cite{komorowski2021minkloclidarmonocularimage}–\cite{LIANG2025408} focus on place recognition, aiming to achieve reliable loop closure detection through multimodal global embeddings. From \cite{Zhao2022} onward, the emphasis shifts toward semantic segmentation, where multimodal fusion enhances pixel- or point-level classification accuracy. Domain adaptation strategies (\cite{Cao2024}, \cite{LuoYang2025}) address distribution shifts between datasets, improving generalization. More advanced tasks include semantic scene completion (\cite{Lu2024}), which infers volumetric occupancy and semantics, and 3D vehicle detection (\cite{Zhao2022}), which strengthens object-level reasoning through feature augmentation. This progression illustrates increasing task complexity and a move from global matching toward dense and structured environment understanding.
The multimodal strategies vary in fusion depth and architectural sophistication. Early approaches (\cite{komorowski2021minkloclidarmonocularimage, Zhao2024}) rely on direct LiDAR–image feature concatenation, while others introduce cross-modal alignment and orientation voting mechanisms (\cite{Xu2025, hou2022pointtovoxelknowledgedistillationlidar, LuoYang2025}) to enforce feature consistency. Attention-based fusion models (\cite{LIANG2025408, Lamichhane2025, Duan2024}) selectively weight complementary modality cues, improving robustness in ambiguous scenes. Transformer-based frameworks (\cite{Zhang2023, Maiti_2023_CVPR, Ma2024}) capture long-range contextual dependencies across modalities, representing a significant architectural advancement. Domain adaptation techniques (\cite{Cao2024,LuoYang2025}) leverage pseudo-label guidance to reduce domain gaps. Task-specific fusion schemes, such as voxelized semantic grids (\cite{Lu2024}) or feature augmentation for detection (\cite{Zhao2022}), tailor multimodal integration to downstream objectives. Overall, there is a clear evolution from simple feature fusion to structured cross-attention and unified encoder architectures.
Across the references, multimodal approaches consistently demonstrate improved robustness and contextual understanding. Methods such as~\cite{komorowski2021minkloclidarmonocularimage} and~\cite{pan2021coralcoloredstructuralrepresentation} effectively integrate geometry and visual structure, enhancing recognition reliability. Cross-modal alignment techniques (\cite{Xu2025, hou2022pointtovoxelknowledgedistillationlidar}) improve orientation consistency and discriminative capacity. Attention-based and transformer-driven models (\cite{LIANG2025408,Zhang2023, Maiti_2023_CVPR,Ma2024}) provide stronger global reasoning and context modeling, leading to improved segmentation accuracy and semantic consistency. Domain adaptation frameworks (\cite{Cao2024,LuoYang2025}) enhance generalization to unseen environments. Scene completion (\cite{Lu2024}) and detection frameworks (\cite{Zhao2022}) benefit from enriched object-aware fusion. In general, multimodal learning leverages complementary sensor characteristics to increase accuracy, resilience, and scene comprehension.
Despite their advantages, these approaches introduce several practical constraints. Many methods (\cite{komorowski2021minkloclidarmonocularimage, Zhou2023, Tan2024}) require accurate LiDAR–camera calibration and synchronization, which may be difficult in real-world deployments. Transformer-based and attention-heavy architectures (\cite{Zhang2023, Maiti_2023_CVPR,Ma2024}) significantly increase computational cost and training complexity. Domain adaptation techniques (\cite{Cao2024,LuoYang2025}) depend on pseudo-label quality and may struggle with severe distribution shifts. Some systems require specialized hardware, such as fisheye cameras (\cite{Zhang2025}), limiting scalability. Additionally, multimodal models often demand large aligned datasets (\cite{Ni2024Robust3S, Zhao2022}) and may exhibit sensitivity to sensor degradation or adverse lighting conditions. Therefore, while multimodal fusion improves perception performance, it also raises challenges in efficiency, scalability, and deployment robustness. \\

In summary, prior work shows that multimodal fusion substantially improves robustness and contextual awareness in both place recognition and dense perception tasks. Nevertheless, many existing approaches either prioritize global retrieval over object-level discrimination or focus on dense prediction without explicitly learning compact and discriminative multimodal descriptors. Moreover, challenges related to efficient fusion, scalability, and generalization remain open. In this context, we propose a method centered on the generation of multimodal fused descriptors that, through learning-based strategies, enable more robust object discrimination and a deeper understanding of the surrounding environment.